\chapter{Instantons as Physical Amplitudes}
\label{ch:instantons-physical-amplitudes}

Chapter~\ref{ch:volume-ii-20} developed the BPST solution, the
finite-regulator instanton density, the anomaly selection rules, and several
QCD amplitude benchmarks.  This chapter reorganizes those ingredients around
the physical object they are meant to compute.  The object is not the
instanton moduli space and not the universal density by itself.  It is a
regulated source or state amplitude, with zero modes saturated, nonzero
fluctuations divided by the chosen counterterms, endpoint faces controlled,
and a Lorentzian or Euclidean observable projection specified.

The distinction is operational.  Two channels may use exactly the same BPST
collective-coordinate space and the same one-loop density, while one channel
vanishes because its sources have rank one in the zero-mode space and the
other gives a nonzero four-fermion coefficient.  Or two Euclidean source
coefficients may agree after integrating over the collective variables, while
their spectral projections differ.  These are not refinements of moduli-space
geometry; they are the steps which turn a saddle contribution into a QFT
amplitude.

\section{The Density Gate: Fluctuations Before Channels}
\label{sec:instanton-density-gate}

The density in \eqref{eq:one-instanton-universal-density} is already a
fluctuation calculation.  The BPST collective-coordinate volume form supplies
one of its factors.  In the half-trace convention of
Chapter~\ref{ch:volume-ii-20}, its universal one-loop part is
\begin{equation}
  \dd\mu_1^{\rm univ}
  =
  C_{\mathcal S,N_c,N_f}\,
  \frac{\dd^4a\,\dd\rho}{\rho^5}\,
  \dd\Omega_{N_c}\,
  \left(\frac{8\pi^2}{g_{\rm ht}^2(\mu)}\right)^{2N_c}
  (\mu\rho)^{b_0}
  \exp\!\left[-\frac{8\pi^2}{g_{\rm ht}^2(\mu)}+\ii\theta\right]
  \left(1+O(g_{\rm ht}^2)\right),
  \label{eq:instanton-physical-density-gate-universal}
\end{equation}
with \(b_0\) defined in
\eqref{eq:instanton-density-beta0-definition}.  The factors have different
origins.  The \(\dd^4a\,\dd\rho\,\rho^{-5}\dd\Omega_{N_c}\) part is the
collective-coordinate Jacobian and orientation convention.  The power
\((8\pi^2/g_{\rm ht}^2)^{2N_c}\) comes from the \(4N_c\) bosonic zero-mode
normalizations.  The exponential is the classical BPST action.  The factor
\((\mu\rho)^{b_0}\) is the logarithmic remnant of the zero-mode-deleted
gauge, ghost, and matter fluctuation determinants after counterterm
subtraction.  The finite constant \(C_{\mathcal S,N_c,N_f}\) is not fixed by
dimensional analysis; it is part of the chosen determinant and orientation
normalization scheme.

\begin{controlledapproximation}[One-loop density gate for a channel]
\label{ca:instanton-one-loop-density-gate-channel}
The one-loop scale test which fixes the logarithmic density power is local:
\[
  \left(\mu\partial_\mu+\beta(g_{\rm ht})\partial_{g_{\rm ht}}\right)
  \log\!\left[
    (\mu\rho)^\alpha
    \exp\!\left(-\frac{8\pi^2}{g_{\rm ht}^2(\mu)}\right)
  \right]
  =
  \alpha-b_0+O(g_{\rm ht}^2).
\]
Thus the density gate passes only for \(\alpha=b_0\).  This test determines
the logarithmic fluctuation power, not the full physical amplitude.  Once the
density is inserted into a channel, the retained integrand can be organized in
the bookkeeping form
\begin{equation}
  \dd\mathcal A_{\mathcal C}^{\rm win}
  =
  \dd\rho\,
  \rho^{b_0-5}\,
  \mathcal W_{\rm fin}^{\rm nz}(\rho,U)\,
  \mathcal Z_{\mathcal C}^{0}(\rho,U)\,
  \mathcal M_{\mathcal C}^{\rm src}(\rho,U)\,
  \mathcal B_{\mathcal C}^{\rho}(\rho)\,
  \dd^4a\,\dd\Omega_{N_c}
  \times(\hbox{scheme and action factors}).
  \label{eq:instanton-density-gate-channel-integrand}
\end{equation}
If the zero-mode/source part starts as
\[
  \mathcal Z_{\mathcal C}^{0}\mathcal M_{\mathcal C}^{\rm src}
  =
  \rho^{\beta_{\mathcal C}}
  \bigl(Z_{\mathcal C}^{(0)}+O(\rho^\delta)\bigr),
\]
then the channel size power is
\(\rho^{b_0+\beta_{\mathcal C}-5}\dd\rho\).  The number
\(\beta_{\mathcal C}\) is physical source data.  It changes when masses
saturate zero modes, when external source derivatives are taken, when
amputation conventions are changed, or when the observable projection kills a
zero-mode component.  For the hard \(SU(3)\), \(N_f=2\), four-source channel
below, \(\beta_{\mathcal C}=6\), so the integrand begins as
\(\rho^{32/3}\dd\rho\), not as the density-only power
\(\rho^{14/3}\dd\rho\).

Consequently a same-theory ratio may cancel the finite determinant constant
and \(\Lambda_{\rm ht}^{b_0}\) only after the source frame, zero-mode
determinants, orientation projection, endpoint window, and physical
projection have been transported.  A computation which keeps
\eqref{eq:instanton-physical-density-gate-universal} but drops
\(\mathcal Z_{\mathcal C}^{0}\), \(\mathcal M_{\mathcal C}^{\rm src}\), or
\(\mathcal B_{\mathcal C}^{\rho}\) has not computed the channel amplitude.
\end{controlledapproximation}

\section{The Channel Is Part of the Calculation}
\label{sec:instanton-physical-amplitude-channel}

The finite-regulator datum
of Definition~\ref{def:one-instanton-amplitude-datum} is the local input for
the present chapter.  We now package it as an observable channel.  This is a
change of emphasis, not a new convention: the BPST/ADHM data select the
domain of integration, while the channel records the determinant, zero-mode,
source, endpoint, and physical-projection data which decide the amplitude.

\begin{definition}[Instanton physical-amplitude channel]
\label{def:instanton-physical-amplitude-channel}
A one-instanton physical-amplitude channel at regulator \(\Lambda\) is a
tuple
\begin{equation}
  \mathcal C_{\Lambda}
  =
  \bigl(
  \mathcal U_{\rm win},
  \dd\nu_{\Lambda,\mathcal M},
  \mathcal W_{\Lambda}^{\rm nz},
  \mathcal Z_{\Lambda}^{0},
  \mathcal M_{\Lambda}^{\rm src},
  \mathcal B_{\Lambda}^{\rho},
  \mathfrak L_{\rm phys}
  \bigr).
  \label{eq:instanton-physical-amplitude-channel}
\end{equation}
Here \(\mathcal U_{\rm win}\) is the retained chart/window in center, size,
and orientation variables; \(\dd\nu_{\Lambda,\mathcal M}\) is the
collective-coordinate density; \(\mathcal W_{\Lambda}^{\rm nz}\) is the
zero-mode-deleted fluctuation determinant with its counterterm convention;
\(\mathcal Z_{\Lambda}^{0}\) is the Berezin coefficient which saturates the
fermion zero modes; \(\mathcal M_{\Lambda}^{\rm src}\) is the finite source,
amputation, OPE, color-singlet, or hadron-overlap map;
\(\mathcal B_{\Lambda}^{\rho}\) is the endpoint, size-window, or screening
factor; and \(\mathfrak L_{\rm phys}\) is the final observable projection,
such as an LSZ residue, spectral-bin functional, Euclidean susceptibility,
OPE coefficient, or inclusive cut.
\end{definition}

The channel coefficient is
\begin{equation}
  \mathcal A_{\Lambda}(\mathcal C)
  =
  \mathfrak L_{\rm phys}
  \left[
  \int_{\mathcal U_{\rm win}}
  \ee^{-S_{\Lambda}(A_z)+\ii\theta}\,
  \mathcal W_{\Lambda}^{\rm nz}(z)\,
  \mathcal Z_{\Lambda}^{0}(z)\,
  \mathcal M_{\Lambda}^{\rm src}(z)\,
  \mathcal B_{\Lambda}^{\rho}(z)\,
  \dd\nu_{\Lambda,\mathcal M}(z)
  \right].
  \label{eq:instanton-physical-channel-functional}
\end{equation}
The physical statement is obtained only with a residual budget:
\begin{equation}
\begin{aligned}
  \mathcal A_{\rm phys}(\mathcal C)
  &=
  \mathcal A_{\Lambda}(\mathcal C)
  +R_{\rm det}+R_{\rm fluc}+R_{\rm zm}+R_{\rm src}+R_{\rm Schur}
   +R_{\rho}                                      \\
  &\qquad
  +R_{\rm sec}+R_{\rm cont}+R_{\rm spec}
   +R_{\rm IR}+R_{\rm cut}+R_{\rm scheme}.
\end{aligned}
  \label{eq:instanton-physical-channel-residual-budget}
\end{equation}
The names match the reusable slots in
Proposition~\ref{prop:instanton-amplitude-assembly-error-budget}.  A
calculation may set some of them to zero only by giving the corresponding
argument; a local BPST density does not set any of the source, endpoint, or
physical-projection terms to zero.

\begin{proposition}[Moduli-equivalent channels need not have the same amplitude]
\label{prop:instanton-moduli-equivalent-channel-separation}
Consider two channels with the same finite instanton chart, the same
collective-coordinate density, and the same nonzero-mode determinant
convention.  In a two-flavor zero-mode sector suppose that the only
channel-dependent local source input is a \(2\times2\) zero-mode source
matrix \(B(z)\), so that
\begin{equation}
  \mathcal Z_{\Lambda}^{0}(z)=
  \det B(z)
  =
  B_{uu}(z)B_{dd}(z)-B_{ud}(z)B_{du}(z).
  \label{eq:instanton-two-flavor-source-determinant}
\end{equation}
Then the common moduli measure and determinant convention do not determine
the amplitude.  If \(B(z)\) has rank one on the retained window, the
two-flavor source contribution vanishes identically, even when the
collective-coordinate density and nonzero-mode determinant are nonzero.  If
another channel has a full-rank \(B'(z)\) on a set of positive retained
weight, its coefficient in
\eqref{eq:instanton-physical-channel-functional} can be nonzero with the same
moduli data.
\end{proposition}

\begin{proof}
At finite regulator, the fermion zero modes are ordinary Grassmann variables.
For two flavors the Berezin integral over the four source variables extracts
the determinant coefficient in
\eqref{eq:instanton-two-flavor-source-determinant}.  A rank-one source matrix
has zero determinant pointwise, so the whole zero-mode coefficient is zero
before the collective-coordinate integral is performed.  A full-rank source
matrix has nonzero determinant at the point under consideration.  Multiplying
by the same bosonic density and the same nonzero-mode determinant therefore
changes the channel coefficient only through the source determinant and the
subsequent observable projection.  The common moduli space has not changed.
\end{proof}

This proposition is the shortest way to state the physics lesson.  The
instanton solution supplies a saddle and a measure factor.  The physical
amplitude asks which external states or sources can see the zero-mode
directions, whether that source information survives projection, and whether
the remaining integral is controlled at its boundaries.

\section{Mass Saturation Versus Source Differentiation}
\label{sec:instanton-mass-source-differentiation}

The local 't Hooft determinant is useful precisely because it keeps the
zero-mode source algebra visible.  It becomes misleading when the determinant
is read as a single undifferentiated number, independent of how the physical
observable is selected.

\begin{proposition}[Two-flavor mass/source determinant coordinate]
\label{prop:two-flavor-source-mass-determinant-coordinate}
Let
\[
  M=\begin{pmatrix}m_u&0\\0&m_d\end{pmatrix},
  \qquad
  B=\begin{pmatrix}B_{uu}&B_{ud}\\B_{du}&B_{dd}\end{pmatrix}
\]
be the mass and source restrictions to the instanton zero-mode space.  The
finite zero-mode Berezin coefficient is
\begin{equation}
  \det(M+B)
  =
  m_um_d
  +m_uB_{dd}+m_dB_{uu}
  +B_{uu}B_{dd}-B_{ud}B_{du}.
  \label{eq:two-flavor-mass-source-determinant-polynomial}
\end{equation}
Thus the mass-saturated vacuum activity, the two-source mass-assisted
coefficients, and the four-source differentiated coefficient are different
coordinates of the same zero-mode determinant.  They are not the same
physical observable.
\end{proposition}

\begin{proof}
This is the determinant of a \(2\times2\) matrix:
\[
  \det(M+B)
  =(m_u+B_{uu})(m_d+B_{dd})-B_{ud}B_{du}.
\]
Expanding gives
\eqref{eq:two-flavor-mass-source-determinant-polynomial}.  Berezin
integration then selects the term with the required number of source
variables or mass insertions.  Setting \(B=0\) leaves the vacuum coordinate
\(m_um_d\); setting \(M=0\) leaves the four-source coefficient
\(\det B\).  Mixed mass/source observables select the linear-in-\(B\)
terms.
\end{proof}

This separation is what the original instanton amplitude calculation uses.
The four-fermion coefficient is obtained by differentiating the source
functional and retaining the zero-mode form factors, not by quoting the
mass-saturated vacuum activity and changing its interpretation.  The
coefficient must then be paired with the nonzero-mode determinant, the color
orientation tensor, external source residues or LSZ data, and a size-window
or infrared mechanism.

\section{A Hard Four-Fermion Benchmark}
\label{sec:instanton-hard-four-fermion-benchmark}

The cleanest test of the preceding distinction is the hard two-flavor
four-fermion channel.  Work in \(SU(3)\) with two light Dirac fundamentals,
choose a \(Q=1\) BPST sector, and insert four amputated external fermion
sources with momenta of size \(Q\).  The center integral first imposes the
momentum-conservation gate.  The color orientation integral then projects the
four source slots through a shared Haar tensor.  The fermion zero modes
produce two rank conditions, one for the right-handed overlap matrix and one
for the left-handed overlap matrix.  Only after these gates have been passed
does the size integral define the hard coefficient.

Let \(Q_\ell=c_\ell Q\), \(c_\ell>0\), be the four hard source scales, and let
\(F_\ell(c_\ell Q\rho)\) denote the individual zero-mode form factor in the
declared source normalization.  After the center delta, common determinant
constant, Haar contraction, amputation convention, and source-overlap
determinants have been pulled into a prefactor \(\mathcal K_{\rm gate}\), the
stripped one-instanton hard coefficient has the form
\begin{equation}
  \mathcal A_{4}^{\rm hard}(Q)
  =
  \mathcal K_{\rm gate}\,
  \Lambda_{\rm ht}^{29/3}\,
  Q^{-35/3}\,
  H(c_1,c_2,c_3,c_4)
  +R_{4}^{\rm hard}(Q),
  \label{eq:instanton-hard-four-source-coefficient}
\end{equation}
where
\begin{equation}
  H(c_1,c_2,c_3,c_4)
  =
  \int_0^\infty
  s^{32/3}\prod_{\ell=1}^{4}F_\ell(c_\ell s)\,\dd s .
  \label{eq:instanton-hard-four-source-window-integral}
\end{equation}
The residual \(R_{4}^{\rm hard}\) contains the determinant, source-frame,
amputation, endpoint, sector-isolation, continuation, and physical-projection
remainders from \eqref{eq:instanton-physical-channel-residual-budget}.  The
formula is a benchmark channel, not a universal local vertex: changing the
source wave packets, the orientation contraction, or the physical projection
changes \(\mathcal K_{\rm gate}\), \(H\), or \(R_{4}^{\rm hard}\).

\begin{proposition}[\(SU(3)\), \(N_f=2\) hard-source power and slow tail]
\label{prop:su3-nf2-hard-source-power-slow-tail}
In the preceding hard channel, \(b_0=29/3\).  The four differentiated
zero-mode slots contribute the local source power \(\rho^6\), so the size
integrand before hard form factors is
\begin{equation}
  \rho^{\,b_0+6-5}\,\dd\rho
  =
  \rho^{32/3}\,\dd\rho .
  \label{eq:instanton-hard-four-source-rho-power}
\end{equation}
After \(s=Q\rho\), the hard coefficient carries the scale power
\begin{equation}
  \Lambda_{\rm ht}^{29/3}Q^{-35/3},
  \qquad
  29/3-35/3=-2 ,
  \label{eq:instanton-hard-four-source-scale-power}
\end{equation}
which is the mass dimension of a four-fermion coefficient.  If all four
individual zero-mode form factors have the large-\(s\) tail
\[
  F_\ell(c_\ell s)=6c_\ell^{-3}s^{-3}+O(s^{-5}),
\]
then the large-\(s\) endpoint of \eqref{eq:instanton-hard-four-source-window-integral}
has leading tail
\begin{equation}
  \int_R^\infty
  s^{32/3}\prod_{\ell=1}^{4}F_\ell(c_\ell s)\,\dd s
  =
  3\,6^4\!\left(\prod_{\ell=1}^{4}c_\ell^{-3}\right)R^{-1/3}
  +O(R^{-7/3}).
  \label{eq:instanton-hard-four-source-slow-tail}
\end{equation}
If one of the four slots is not hard, this power test no longer controls the
large-size endpoint by the same source profile.
\end{proposition}

\begin{proof}
For \(SU(3)\) with two Dirac fundamentals,
\[
  b_0=\frac{11}{3}3-\frac{2}{3}2=\frac{29}{3}.
\]
The universal one-loop size density contributes \(\rho^{b_0-5}\dd\rho\), and
the four differentiated zero-mode slots contribute \(\rho^6\).  This gives
\eqref{eq:instanton-hard-four-source-rho-power}.  The change of variables
\(s=Q\rho\) gives
\[
  \dd\rho\,\rho^{32/3}
  =
  Q^{-35/3}s^{32/3}\,\dd s .
\]
At the hard subtraction point \(\mu=Q\), the one-loop density prefactor
\(\mu^{b_0}\) contributes \(Q^{29/3}\), so before rewriting the classical
exponential the coefficient scales as
\(\exp[-8\pi^2/g_{\rm ht}^2(Q)]Q^{-2}\).  The one-loop RG-invariant trade
\[
  \exp[-8\pi^2/g_{\rm ht}^2(Q)]
  =
  \Lambda_{\rm ht}^{29/3}Q^{-29/3}
\]
then gives the final power
\(\Lambda_{\rm ht}^{29/3}Q^{-35/3}\).  The total mass dimension is
\(29/3-35/3=-2\), as required for a coefficient multiplying a dimension-six
four-fermion source.

For the endpoint, multiply the four large-\(s\) tails.  The product contributes
\[
  6^4\left(\prod_\ell c_\ell^{-3}\right)s^{-12}+O(s^{-14}).
\]
Multiplication by \(s^{32/3}\) gives
\[
  6^4\left(\prod_\ell c_\ell^{-3}\right)s^{-4/3}+O(s^{-10/3}).
\]
Integrating from \(R\) to infinity gives
\eqref{eq:instanton-hard-four-source-slow-tail}.  With only three hard slots
the leading endpoint power is \(s^{32/3-9}=s^{5/3}\), which is not integrable
at infinity.
\end{proof}

\begin{controlledapproximation}[Hard benchmark gate ledger]
\label{ca:instanton-hard-benchmark-gate-ledger}
A controlled comparison of two hard four-source instanton coefficients at
scales \(Q_1,Q_2\) may use the power in
\eqref{eq:instanton-hard-four-source-scale-power} only after the gate data are
held fixed:
\[
  \Delta_{\rm cen}\ne0,\qquad
  \mathcal P_{\rm Haar}\ne0,\qquad
  \det B_R\ne0,\qquad
  \det B_L\ne0,
\]
with the same amputation convention, source-frame normalization, determinant
scheme, sector projection, and dimensionless source-window shape.  If the
retained coefficients have the form
\[
  \mathcal A_j
  =
  \mathcal A_j^{(0)}(1+\varepsilon_j),
  \qquad
  |\varepsilon_j|\le e_j<1,
  \qquad j=1,2,
\]
and the leading gate data are the same at the two scales, then
\begin{equation}
  \frac{\mathcal A_2}{\mathcal A_1}
  =
  \left(\frac{Q_2}{Q_1}\right)^{-35/3}
  \frac{1+\varepsilon_2}{1+\varepsilon_1},
  \label{eq:instanton-hard-four-source-ratio}
\end{equation}
with residual multiplier bounded by
\begin{equation}
  \left|
  \frac{1+\varepsilon_2}{1+\varepsilon_1}-1
  \right|
  \le
  \frac{e_1+e_2}{1-e_1}.
  \label{eq:instanton-hard-four-source-ratio-bound}
\end{equation}
Thus the determinant constant and \(\Lambda_{\rm ht}\) cancel in a same-theory
hard-scale ratio, while the dimensionless source-window and rank-conditioning
data do not cancel unless they have actually been transported.  An off-shell
center gate, a rank-one zero-mode overlap, an omitted Haar projection, or
unamputated external residues change the amplitude coordinate before the power
law \((Q_2/Q_1)^{-35/3}\) can be tested.
\end{controlledapproximation}

\section{Finite-Cell Control Model}
\label{sec:finite-cell-instanton-channel-control}

A continuum proof of an instanton amplitude is analytic.  Nevertheless, the
logic which prevents a moduli-only shortcut is already visible in a finite
cell model.  Let the retained instanton window be decomposed into finitely
many cells \(i\).  Let
\[
  c_i
  =
  w_i\,d_i\,\det B_i\,p_i
\]
denote the leading cell contribution after the positive collective weight
\(w_i\), determinant factor \(d_i\), zero-mode source determinant
\(\det B_i\), and physical projection weight \(p_i\) have all been inserted.

\begin{controlledapproximation}[Finite-cell instanton channel control]
\label{ca:finite-cell-instanton-channel-control}
In a finite retained-cell channel, suppose each leading cell contribution
\(c_i\) is replaced by \(c_i(1+\delta_i)\), with
\(|\delta_i|\le\epsilon_i\), and suppose the omitted projection, sector,
infrared, cut, and scheme remainders have absolute majorant \(R_{\rm ext}\).
Then
\begin{equation}
  \left|
  \mathcal A_{\rm phys}
  -\sum_i c_i
  \right|
  \le
  \sum_i |c_i|\,\epsilon_i+R_{\rm ext}.
  \label{eq:finite-cell-instanton-channel-bound}
\end{equation}
If the source matrix on a cell is perturbed by
\(B_i\mapsto B_i+E_i\), with \(B_i\) invertible and
\(\eta_i=\|B_i^{-1}E_i\|\) in an operator norm for which
\(|\operatorname{tr}X|\le2\|X\|\) and \(|\det X|\le\|X\|^2\) on
\(2\times2\) matrices, then
\begin{equation}
  \frac{|\det(B_i+E_i)-\det B_i|}{|\det B_i|}
  \le
  2\eta_i+\eta_i^2 .
  \label{eq:finite-cell-source-determinant-stability}
\end{equation}
The first estimate is the triangle inequality applied after the channel has
already been projected to the same observable coordinate.  It would be false
as a relative statement without a noncancellation margin for
\(\sum_i c_i\).  For the determinant estimate, write
\[
  \det(B_i+E_i)=\det B_i\,
  \det(1+B_i^{-1}E_i).
\]
For a \(2\times2\) matrix \(X\),
\(\det(1+X)-1=\operatorname{tr}X+\det X\).  The stated norm bounds give
\eqref{eq:finite-cell-source-determinant-stability}.
\end{controlledapproximation}

The finite-cell model has two uses.  First, it gives a transparent negative
control: deleting \(\det B_i\), \(p_i\), or \(R_{\rm ext}\) changes the
physical statement, even if the collective-coordinate weights are unchanged.
Second, it records why the continuum estimates in
Chapter~\ref{ch:volume-ii-20} carry so many named residuals.  They
are not decorative error terms.  They are the places where a Euclidean saddle
coefficient becomes, or fails to become, a physical amplitude.
